\section{Organic Chemistry}

\subsection{Hydrocarbons}

\subsubsection{Alkanes}
\[
\mathrm{C}_n\mathrm{H}_{2n+2}
\]
Alkanes are saturated hydrocarbons with only single bonds. They are named by the longest carbon chain and alkyl substituents.

\subsubsection{Alkenes}
\[
\mathrm{C}_n\mathrm{H}_{2n}
\]
Alkenes contain a carbon-carbon double bond. Restricted rotation can give cis/trans isomerism when each alkene carbon has different substituents.

\subsubsection{Alkynes}
\[
\mathrm{C}_n\mathrm{H}_{2n-2}
\]
Alkynes contain a carbon-carbon triple bond. \(\ce{C2H2}\) is the simplest alkyne Lewis-structure example in the slides.

\subsubsection{Aromatic hydrocarbons}
\[
\ce{C6H6}
\]
Benzene is a planar six-membered aromatic ring with one hydrogen on each carbon. Aromatic rings are recurring structural motifs in functional-group questions.

\subsection{Isomerism and Nomenclature}

\subsubsection{Structural isomers}
\[
\text{same molecular formula, different atom connectivity}
\]
Isomers have different structures and therefore different properties. Alkane isomer counting appears in lecture exercises.

\subsubsection{Alkane naming}
\[
\text{locants}+\text{substituents}+\text{parent chain}
\]
Choose the longest carbon chain, number from the end nearest substituents, and list substituents alphabetically.

\subsubsection{Degree of unsaturation}
\[
\mathrm{DBE}=\frac{2C+2+N-H-X}{2}
\]
The double-bond equivalent counts rings and \(\pi\)-bonds for formulas containing C, H, N, and halogens \(X\). Oxygen does not affect the count.

\subsection{Functional Groups}

\subsubsection{Functional group recognition}
\[
\begin{array}{c c}
\toprule
\text{group} & \text{structural motif}\\
\midrule
\text{alcohol} & \mathrm{R{-}OH}\\
\text{ether} & \mathrm{R{-}O{-}R'}\\
\text{aldehyde} & \mathrm{R{-}CHO}\\
\text{ketone} & \mathrm{R{-}CO{-}R'}\\
\text{carboxylic acid} & \mathrm{R{-}COOH}\\
\text{ester} & \mathrm{R{-}COO{-}R'}\\
\text{amine} & \mathrm{R{-}NH_2},\mathrm{R_2NH},\mathrm{R_3N}\\
\text{amide} & \mathrm{R{-}CONH_2},\mathrm{R{-}CONHR'},\mathrm{R{-}CONR'_2}\\
\bottomrule
\end{array}
\]
Functional group identification is a recurring exam concept, especially in molecules with several groups.

\subsubsection{Carboxylic acid deprotonation}
\[
\mathrm{R{-}COOH}\rightleftharpoons \mathrm{R{-}COO^-}+\ce{H+}
\]
Carboxylic acids contain a \(\ce{COOH}\) group and form carboxylate conjugate bases.

\subsubsection{Ester and ether distinction}
\[
\text{ester: }\mathrm{R{-}C(=O){-}O{-}R'}, \qquad
\text{ether: }\mathrm{R{-}O{-}R'}
\]
An ester contains a carbonyl next to oxygen. An ether contains an oxygen single-bonded to two carbon groups.

\subsection{Chirality and Biomolecules}

\subsubsection{Chiral center}
\[
\mathrm{C^*}=\text{tetrahedral carbon attached to four different groups}
\]
Interchanging two groups at a chiral center gives the enantiomer. Chirality was highlighted with ibuprofen and amino-acid structures.

\subsubsection{Amino acids and peptides}
\[
\mathrm{H_2N{-}CHR{-}COOH}
\]
Amino acids contain amine and carboxylic acid functional groups. Peptides form by connecting amino acids through amide bonds.

\subsubsection{Degree of polymerization}
\[
n=\frac{M_{\mathrm{polymer}}}{M_{\mathrm{repeat}}}
\]
% Added from exam coverage
The degree of polymerization is the number of repeat units in a polymer chain.
